\chapter{Introduction}
A communications network is a system made up of interconnected devices that exchange information through a set of rules. In this context, the dominant transport paradigm is packet switching, where data is divided into independent units (packets) that are transmitted across the network without a predefined fixed route.

The organization of these systems is defined by network architectures, which describe both the physical elements (nodes and links) and the logical elements (protocols, addressing, and communication rules). The most widely adopted architecture is TCP/IP (Transmission Control Protocol / Internet Protocol), originally designed for the interconnection of heterogeneous networks — a concept from which the term Internet derives (inter-networking).

TCP/IP organizes communication into layers, facilitating interoperability between systems. This approach was later formalized by the theoretical OSI model (Open Systems Interconnection), which defines seven layers (physical → data link → network → transport → session → presentation → application), although in practice the Internet operates on a simplified four-layer version (physical → network → transport → application). Each layer groups different protocols that implement specific functions.

Within this architecture, one of the fundamental protocols is DNS (Domain Name System), located at the application layer. Its function is to translate human-readable domain names (e.g., amazon.com) into the IP addresses used by the network. DNS is a hierarchical and distributed system whose operation involves a sequence of queries: the client contacts a recursive resolver (Google, Cloudflare, etc.), which in turn interacts with root servers, TLD (Top-Level Domain) servers, and finally authoritative servers to obtain the requested IP address.
However, the original design of DNS does not include any security mechanisms such as encryption or authentication. Queries and responses are transmitted in plain text, exposing sensitive information about user activity. This lack of privacy makes DNS traffic a relevant target for various types of attacks.

These attacks can be grouped into three main categories. First, attacks on resolution integrity, such as cache poisoning (storing malicious resolutions in the memory of resolver nodes), DNS spoofing (intercepting and modifying DNS traffic), or hijacking (taking control of a resolver node and therefore its responses) — all of which aim to manipulate DNS responses. Second, DDoS (Distributed Denial of Service) attacks, such as amplification or flood attacks, which seek to overwhelm the infrastructure. Finally, there is protocol abuse through techniques such as DNS tunneling, where DNS is used as a covert channel for data exfiltration or Command and Control (C2) operations.

To mitigate confidentiality and integrity issues, solutions based on the use of security protocols have been introduced. A key precedent is HTTPS, which combines HTTP (HyperText Transfer Protocol) with TLS (Transport Layer Security), providing encryption, integrity, and authentication. Following this same approach, secure variants of DNS have emerged — DoT (DNS over TLS), DoH (DNS over HTTPS), and later DoQ (DNS over QUIC) — all of which encapsulate DNS queries within encrypted channels.

These solutions protect communication between the client and the resolver, reducing man-in-the-middle (MITM) attacks and passive traffic monitoring. However, they do not eliminate all attack vectors; rather, they transform the problem. In particular, they hinder the inspection of DNS content, which reduces the effectiveness of traditional detection techniques based on DPI (Deep Packet Inspection) or heuristic rules.

As a result, threat detection in encrypted DNS traffic shifts toward the analysis of metadata and behavioral patterns. In this context, ML (Machine Learning)-based approaches have emerged, evolving from more classical paradigms such as SL (Statistical Learning) to more modern ones such as DL (Deep Learning). Given the dynamic nature of threats, paradigms such as CL (Continuous Learning) have been incorporated, and more recently FL (Federated Learning), which enables distributed model training without sharing sensitive data — aligning with the privacy requirements introduced by these new secure DNS technologies.

This paradigm shift gives rise to the central question of this work: to assess to what extent these methods are capable of detecting threats in environments where traffic visibility is limited by encryption.
